\documentclass[11pt,a4paper]{article}

\usepackage[T1]{fontenc}
\usepackage[utf8]{inputenc}
\usepackage[a4paper,margin=27mm]{geometry}
\usepackage{amsmath,amssymb,amsthm,mathtools}
\usepackage{newpxtext,newpxmath}
\usepackage{microtype}
\usepackage{booktabs,tabularx,array}
\usepackage{enumitem}
\usepackage{xcolor}
\usepackage{titlesec}
\usepackage{mdframed}
\usepackage[authoryear,round]{natbib}
\usepackage[colorlinks=true,linkcolor=blue!45!black,
  citecolor=red!45!black,urlcolor=blue!45!black]{hyperref}

\definecolor{otblue}{HTML}{1F4E79}
\definecolor{otgray}{HTML}{F1F3F5}
\definecolor{otline}{HTML}{AAB5BF}
\definecolor{otremark}{HTML}{6D8E9F}

\titleformat{\section}
  {\normalfont\Large\bfseries\color{otblue}}
  {\thesection}{.7em}{}
\titleformat{\subsection}
  {\normalfont\large\bfseries\color{otblue!80!black}}
  {\thesubsection}{.65em}{}
\titleformat{\paragraph}[runin]
  {\normalfont\bfseries\color{otblue!75!black}}
  {}{0pt}{}
\titlespacing*{\section}{0pt}{2.7ex plus .7ex minus .2ex}{1.1ex}
\titlespacing*{\subsection}{0pt}{2.0ex plus .5ex minus .2ex}{.75ex}
\titlespacing*{\paragraph}{0pt}{1.3ex plus .3ex minus .1ex}{.55em}

\newtheorem{theorem}{Theorem}[section]
\newtheorem{proposition}[theorem]{Proposition}
\newtheorem{lemma}[theorem]{Lemma}
\newtheorem{corollary}[theorem]{Corollary}
\theoremstyle{definition}
\newtheorem{example}[theorem]{Example}
\newtheorem{remark}[theorem]{Remark}

\newmdenv[
  backgroundcolor=otgray,
  linecolor=otline,
  linewidth=.5pt,
  roundcorner=2pt,
  innerleftmargin=8pt,
  innerrightmargin=8pt,
  innertopmargin=6pt,
  innerbottommargin=6pt,
  skipabove=8pt,
  skipbelow=8pt
]{definitionbox}

\newenvironment{definition}[1][]{%
  \begin{definitionbox}\refstepcounter{theorem}%
  \noindent\textbf{Definition~\thetheorem}%
  \if\relax\detokenize{#1}\relax.\else\ \textbf{(#1).}\fi\enspace
}{\end{definitionbox}}

\newcommand{\RR}{\mathbb R}
\newcommand{\Pp}{\mathcal P}
\newcommand{\Ptwo}{\mathcal P_2}
\newcommand{\Wass}{\mathcal W}
\newcommand{\Wgrad}{\nabla_{\!\Wass}}
\newcommand{\Id}{\operatorname{Id}}
\newcommand{\argmin}{\operatorname*{arg\,min}}
\newcommand{\diverg}{\operatorname{div}}
\newcommand{\supp}{\operatorname{supp}}
\newcommand{\KL}{\operatorname{KL}}
\newcommand{\Ent}{\operatorname{Ent}}
\newcommand{\Law}{\operatorname{Law}}
\newcommand{\dd}{\,\mathrm d}
\newcommand{\eps}{\varepsilon}
\newcommand{\norm}[1]{\left\lVert#1\right\rVert}
\newcommand{\abs}[1]{\left\lvert#1\right\rvert}
\newcommand{\pair}[2]{\left\langle#1,#2\right\rangle}
\newcommand{\al}{\alpha}
\newcommand{\be}{\beta}
\newcommand{\R}{\mathbb R}
\newcolumntype{P}[1]{>{\raggedright\arraybackslash}p{#1}}
\newcolumntype{Y}{>{\raggedright\arraybackslash}X}

\title{A Gentle Introduction to Mean Field Games\\
  \large Optimal transport, variational problems on measures, and forward--backward PDEs}
\author{Gabriel Peyr\'e}
\date{July 2026}

\hypersetup{
  pdftitle={A Gentle Introduction to Mean Field Games},
  pdfauthor={Gabriel Peyre},
  pdfsubject={Mean field games, optimal transport, and forward-backward PDEs},
  pdfkeywords={mean field games, optimal transport, JKO scheme, Hamilton-Jacobi-Bellman equation, Fokker-Planck equation}
}

\begin{document}
\maketitle

\begin{abstract}
Mean field games describe Nash equilibria in populations containing very many
weakly interacting decision makers.  Their basic mathematical object is a
forward--backward PDE system: a Hamilton--Jacobi--Bellman equation propagates
the value of a representative agent backward in time, while a continuity or
Fokker--Planck equation propagates the population law forward.  This note gives
a gentle introduction using the notation and viewpoint of \emph{Optimal
Transport for Machine Learners}.  The emphasis is on the links with transport
of probability measures, Benamou--Brenier actions, first variations of
functionals, potential games, Wasserstein geometry, and minimizing movements.
It also stresses an important boundary: a generic mean field game is a Nash
equilibrium problem, not a Wasserstein gradient flow.  The JKO scheme and the
global variational formulation of a potential game are closely related but
encode different time structures.  After deriving the MFG system, we discuss
monotonicity, uniqueness, finite-player limits, master equations, examples,
and the main numerical strategies.
\end{abstract}

\tableofcontents

\section{Why mean field games?}

Mean field games (MFGs) were introduced independently by
\citet{LasryLions2006a,LasryLions2006b,LasryLions2007} and by
\citet{HuangMalhameCaines2006} to describe differential games with a very
large number of statistically similar agents.  The motivating compromise is
simple.  A game with $N$ players has an $N$-body state and, in feedback
form, an $N$-body system of Hamilton--Jacobi equations.  Direct computation
quickly becomes impossible.  Yet replacing the players by one centrally
controlled fluid changes the model: selfish agents seek a Nash equilibrium,
not the minimizer of a social cost.  The mean field limit keeps the strategic
character of the game while compressing the other $N-1$ players into their
empirical distribution.

The resulting state variable is a probability measure
$\al_t\in\Pp(\mathcal X)$.  This makes MFGs naturally close to optimal
transport, mean-field particle systems, and Wasserstein gradient flows.  The
same continuity and Fokker--Planck equations occur in all these theories, but
their interpretations differ:

\begin{itemize}[leftmargin=2em]
  \item in optimal transport, a velocity is chosen to move one prescribed
  measure to another with minimal kinetic action;
  \item in a Wasserstein gradient flow, the velocity follows the steepest
  descent of one functional of the current measure;
  \item in mean field control, a planner chooses the population dynamics to
  minimize one collective objective;
  \item in a mean field game, every infinitesimal agent optimizes against a
  population curve that must, at equilibrium, be generated by the agents'
  optimal responses.
\end{itemize}

This last fixed-point requirement is the source of the characteristic
forward--backward coupling.  The population law determines the value function;
the value function determines an optimal feedback; and that feedback must
reproduce the population law.

\paragraph{Scope.}
We work mostly on the flat torus $\mathcal X=\mathbb T^d$, which avoids
boundary conditions, or on $\RR^d$ under sufficient moment and decay
assumptions.  We present smooth equations first and indicate where weak or
viscosity solutions are needed.  Systematic treatments include
\citet{CardaliaguetNotes2010,BensoussanFrehseYam2013,
CarmonaDelarue2018a,CarmonaDelarue2018b,CardaliaguetDelarueLasryLions2019};
the numerical viewpoint is surveyed by \citet{AchdouLauriere2020}.

\section{From an \texorpdfstring{$N$}{N}-player game to a population law}

\paragraph{The finite game.}
Consider $N$ exchangeable players.  Player $i$ has state
$X_t^i\in\RR^d$ and chooses a progressively measurable control
$a_t^i\in\RR^d$.  In the simplest model,
\begin{equation}\label{eq:n-player-dynamics}
  \dd X_t^i=a_t^i\,\dd t+\sqrt{2\nu}\,\dd B_t^i,
  \qquad \nu\geq0,
\end{equation}
where the Brownian motions are independent.  The other players enter only
through the empirical measure
\begin{equation}\label{eq:empirical-population}
  \al_t^{N,-i}=\frac1{N-1}\sum_{j\neq i}\delta_{X_t^j}.
\end{equation}
Given a running Lagrangian $L(x,a)$, a running interaction $F(x,\al)$, and
a terminal interaction $G(x,\al)$, player $i$ minimizes
\begin{equation}\label{eq:n-player-cost}
  J_i^N(a^1,\ldots,a^N)
  =\mathbb E\!\left[
    \int_0^T \bigl(L(X_t^i,a_t^i)+F(X_t^i,\al_t^{N,-i})\bigr)\,\dd t
    +G(X_T^i,\al_T^{N,-i})
  \right].
\end{equation}

\begin{definition}[Nash equilibrium]
A control profile $(a^{1,*},\ldots,a^{N,*})$ is a Nash equilibrium if no
player can lower its own cost by changing only its control: for every $i$
and every admissible $a^i$,
\[
  J_i^N(a^{1,*},\ldots,a^{N,*})
  \leq
  J_i^N(a^{1,*},\ldots,a^{i-1,*},a^i,a^{i+1,*},\ldots,a^{N,*}).
\]
It is an $\eps$-Nash equilibrium if the right-hand side may improve the cost
by at most $\eps$.
\end{definition}

The equilibrium is non-cooperative: the feedback of player $i$ minimizes
its own objective while the strategies of all other players are frozen.  This
is different from minimizing $\sum_iJ_i^N$, which is a centralized control
problem.

\paragraph{The mean field ansatz.}
When $N$ is large and each player has weight $1/N$, one expects a typical
agent to have a negligible effect on the population law.  The empirical
measure is replaced by a deterministic curve
$(\al_t)_{0\leq t\leq T}$, and the representative agent solves a standard
control problem with this curve frozen.  The optimal controlled state must
then have law $\al_t$.  Symbolically,
\[
  \text{population curve}
  \longmapsto
  \text{optimal feedback}
  \longmapsto
  \text{law of the optimally controlled state}.
\]
An MFG equilibrium is a fixed point of this map.

Under suitable regularity and uniqueness assumptions, the feedback obtained
from the MFG system yields an $\eps_N$-Nash equilibrium for the finite game,
with $\eps_N\to0$.  The probabilistic mechanism is propagation of chaos:
finitely many tagged players become asymptotically independent with common law
$\al_t$.  Precise open-loop and closed-loop results require different
arguments; see \citet{CarmonaDelarueLacker2016,Lacker2020,
CardaliaguetDelarueLasryLions2019}.

\section{The representative-agent control problem}

Fix for the moment a continuous curve
$t\mapsto\al_t\in\Ptwo(\RR^d)$.  Starting from $X_t=x$, a representative
agent chooses $a_s$ in
\begin{equation}\label{eq:representative-sde}
  \dd X_s=a_s\,\dd s+\sqrt{2\nu}\,\dd B_s,
  \qquad s\in[t,T],
\end{equation}
and minimizes
\begin{equation}\label{eq:representative-cost}
  \mathbb E\!\left[
    \int_t^T\bigl(L(X_s,a_s)+F(X_s,\al_s)\bigr)\,\dd s
    +G(X_T,\al_T)
  \right].
\end{equation}
Its value is denoted by $u(t,x)$.

We use the Hamiltonian convention
\begin{equation}\label{eq:mfg-hamiltonian}
  H(x,p)=\sup_{a\in\RR^d}\bigl\{-a\cdot p-L(x,a)\bigr\}.
\end{equation}
If $L(x,\cdot)$ is strictly convex and smooth, the maximizing control is
\begin{equation}\label{eq:optimal-feedback}
  a^*(x,p)=-\nabla_pH(x,p).
\end{equation}
For $L(x,a)=\abs a^2/2$, one has $H(x,p)=\abs p^2/2$ and $a^*=-p$.

\begin{proposition}[Hamilton--Jacobi--Bellman equation]
\label{prop:hjb-fixed-population}
Assume temporarily that the value function is smooth.  For a prescribed
population curve $(\al_t)_t$, the value function of
\eqref{eq:representative-cost} satisfies
\begin{equation}\label{eq:mfg-hjb-fixed}
  -\partial_tu-\nu\Delta u+H(x,\nabla u)=F(x,\al_t),
  \qquad
  u(T,x)=G(x,\al_T).
\end{equation}
The optimal Markov feedback is
$a_t^*(x)=-\nabla_pH(x,\nabla u(t,x))$.
\end{proposition}

\begin{proof}
The dynamic-programming principle over a short interval $[t,t+h]$ gives
\[
u(t,x)=\inf_a\mathbb E\!\left[
  \int_t^{t+h}(L(X_s,a_s)+F(X_s,\al_s))\,\dd s
  +u(t+h,X_{t+h})\right].
\]
Apply It\^o's formula to $u(s,X_s)$, divide by $h$, and let $h\downarrow0$.
The local minimization is
$\inf_a\{L(x,a)+a\cdot\nabla u\}=-H(x,\nabla u)$, which gives
\eqref{eq:mfg-hjb-fixed}.  The optimizer is
\eqref{eq:optimal-feedback}.
\end{proof}

When smoothness fails, the HJB equation is naturally interpreted in the
viscosity sense.  This is one reason why the backward equation and the forward
equation do not live in identical function spaces.

\section{The mean field game PDE system}

\paragraph{Closing the fixed point.}
Apply the optimal feedback from Proposition~\ref{prop:hjb-fixed-population} to
an initial random state $X_0\sim\al_0$.  Its law solves the Kolmogorov forward
equation.  Requiring this law to equal the curve used in the control problem
closes the system.

\begin{definition}[Finite-horizon mean field game system]
Given an initial law $\al_0$, a classical MFG equilibrium is a pair
$(u,\al)$ satisfying
\begin{equation}\label{eq:mfg-system}
\left\{
\begin{aligned}
 -\partial_tu_t-\nu\Delta u_t+H(x,\nabla u_t)
   &=F(x,\al_t),
   &&u_T(x)=G(x,\al_T),\\
 \partial_t\al_t-\nu\Delta\al_t
   -\diverg\!\left(\al_t\nabla_pH(x,\nabla u_t)\right)
   &=0,
   &&\al_{t=0}=\al_0.
\end{aligned}
\right.
\end{equation}
The second equation is understood distributionally when $\al_t$ has no
density.
\end{definition}

The first equation runs backward because its terminal value is prescribed;
the second runs forward because its initial law is prescribed.  Neither can be
solved independently: $F$ and $G$ couple the population into the HJB
equation, while $u$ determines the drift in the Fokker--Planck equation.

\begin{proposition}[Weak forward equation and mass conservation]
\label{prop:mfg-weak-fp}
Let $b_t(x)=-\nabla_pH(x,\nabla u_t(x))$.  The forward equation in
\eqref{eq:mfg-system} is equivalent to
\begin{equation}\label{eq:mfg-weak-fp}
  \frac{\dd}{\dd t}\int\varphi(x)\,\dd\al_t(x)
  =\int\bigl(\nabla\varphi(x)\cdot b_t(x)
      +\nu\Delta\varphi(x)\bigr)\,\dd\al_t(x)
\end{equation}
for every smooth test function $\varphi$.  In particular,
$\al_t(\mathcal X)=\al_0(\mathcal X)$.
\end{proposition}

\begin{proof}
Multiply the Fokker--Planck equation by $\varphi$ and integrate by parts.
Taking $\varphi\equiv1$ gives mass conservation.
\end{proof}

For $\nu=0$, the forward equation is exactly the continuity equation
\begin{equation}\label{eq:mfg-continuity}
  \partial_t\al_t+\diverg(\al_tv_t)=0,
  \qquad
  v_t=-\nabla_pH(x,\nabla u_t).
\end{equation}
This is the same Eulerian kinematics used by dynamic optimal transport and
Wasserstein gradient flows.  For $\nu>0$, it is a Fokker--Planck equation.
If $\al_t=\rho_t\dd x$ is smooth and positive, it can also be rewritten as a
continuity equation with current velocity
\begin{equation}\label{eq:mfg-current-velocity}
  \partial_t\rho_t+\diverg(\rho_t\widetilde v_t)=0,
  \qquad
  \widetilde v_t=-\nabla_pH(x,\nabla u_t)-\nu\nabla\log\rho_t.
\end{equation}
The score correction is the same one that relates Langevin diffusion and
deterministic probability-flow representations.

\begin{remark}[What is forward and what is backward]
The density equation is causal once a feedback is known.  The feedback is not
causal in the population curve because a rational agent anticipates future
congestion and the terminal cost.  This anticipation is encoded by the
backward HJB equation.  Setting $u_t$ from an instantaneous cost only would
produce a myopic interacting-particle or gradient-flow model, not the same
game.
\end{remark}

\section{An optimal-transport dictionary}

The analogy with the rest of the book is clearest at the level of mathematical
objects rather than slogans.

\begin{center}
\begin{tabularx}{\textwidth}{>{\bfseries}P{.24\textwidth}Y}
\toprule
Optimal-transport/PDE object & Role in a mean field game\\
\midrule
Probability measure $\al_t$ & Distribution of the state of a representative agent.\\
Empirical measure & Finite-player approximation of the population law.\\
Continuity equation & Deterministic transport of the population by the equilibrium feedback.\\
Fokker--Planck equation & Transport plus idiosyncratic Brownian exploration.\\
Kinetic action & Individual effort, and in potential games the dynamic transport term of the social potential.\\
Hamilton--Jacobi equation & Dual/optimality equation for a representative agent.\\
First variation & Converts a potential functional of $\al$ into the coupling cost perceived by one agent.\\
JKO step & A proximal, myopic evolution of measures; related to but distinct from a finite-horizon MFG equilibrium.\\
Geodesic convexity & A useful convexity mechanism for potential planning problems, not a replacement for Nash monotonicity.\\
\bottomrule
\end{tabularx}
\end{center}

\paragraph{The Benamou--Brenier limit.}
If $F=0$, $G$ imposes a hard terminal law, $\nu=0$, and
$L(a)=\abs a^2/2$, the population planning problem reduces to dynamic
quadratic optimal transport.  The HJB equation becomes the Hamilton--Jacobi
equation dual to the Benamou--Brenier formulation, and the optimal population
velocity is a gradient.  MFGs enlarge this picture by adding endogenous
running and terminal interactions, while generic games lose the global
variational structure entirely \citep{BenamouBrenier2000}.

\paragraph{Schr\"odinger bridges.}
With $\nu>0$ and prescribed endpoint laws, minimizing quadratic drift energy
relative to Brownian motion gives the Schr\"odinger bridge.  A second-order
potential MFG has a similar controlled-diffusion constraint, but its final law
is usually endogenous and its congestion or interaction potential is
integrated along the path.  The entropy-on-path-space representation of
variational second-order MFGs is developed by
\citet{BenamouCarlierDiMarinoNenna2019}.

\section{Potential mean field games as variational problems}

A generic MFG is only a fixed point.  A particularly important subclass is
generated by functionals on the space of measures.

\begin{definition}[Potential couplings]
The running and terminal couplings are potential if there exist functionals
$\mathcal F,\mathcal G:\Pp(\mathcal X)\to\RR\cup\{+\infty\}$ such that
\begin{equation}\label{eq:potential-first-variations}
  \frac{\delta\mathcal F}{\delta\al}(\al)(x)=F(x,\al),
  \qquad
  \frac{\delta\mathcal G}{\delta\al}(\al)(x)=G(x,\al),
\end{equation}
up to additive constants independent of $x$.
\end{definition}

For clarity, first take $L(a)=\abs a^2/2$.  The potential planning problem is
\begin{equation}\label{eq:potential-mfg-primal}
\inf_{(\al,v)}
\left\{
  \int_0^T\left[
    \frac12\int\abs{v_t(x)}^2\,\dd\al_t(x)+\mathcal F(\al_t)
  \right]\dd t
  +\mathcal G(\al_T)
\right\},
\end{equation}
subject to
\begin{equation}\label{eq:potential-mfg-constraint}
  \partial_t\al_t+\diverg(\al_tv_t)=\nu\Delta\al_t,
  \qquad \al_{t=0}=\al_0.
\end{equation}
The unknown final law is selected by the terminal potential.  Prescribing
$\al_T=\al_1$ instead gives the mean field planning problem.

\begin{proposition}[Optimality system of a quadratic potential MFG]
\label{prop:potential-mfg-optimality}
Assume all fields are smooth, $\al_t=\rho_t\dd x$ is positive, and
\eqref{eq:potential-mfg-primal} admits a minimizer.  Then its first-order
optimality conditions are
\begin{equation}\label{eq:potential-mfg-system}
\left\{
\begin{aligned}
 -\partial_tu-\nu\Delta u+\frac12\abs{\nabla u}^2
   &=F(x,\al_t),\\
 \partial_t\al_t-\nu\Delta\al_t-\diverg(\al_t\nabla u)&=0,\\
 \al_{t=0}&=\al_0,
 \qquad u_T=G(\cdot,\al_T).
\end{aligned}
\right.
\end{equation}
Thus a minimizer gives an MFG equilibrium.
\end{proposition}

\begin{proof}
Introduce the multiplier $-u$ for
\eqref{eq:potential-mfg-constraint}, namely add
\[
 -\int_0^T\!\int u
  (\partial_t\rho+\diverg(\rho v)-\nu\Delta\rho)\,\dd x\dd t
\]
to the objective.
After integration by parts, the terms depending on $v$ are
$\int\rho(\abs v^2/2+v\cdot\nabla u)$.  Their minimizer is
$v=-\nabla u$.  Variation in $\rho$, together with
\eqref{eq:potential-first-variations}, gives the HJB equation.  The free
terminal variation gives $u_T=G(\cdot,\al_T)$.  These are precisely
\eqref{eq:potential-mfg-system}.
\end{proof}

The sign bookkeeping in the proof is conventional; the invariant statement
is that minimizing the kinetic term produces the Legendre relation between
the velocity and the HJB momentum.  For a general convex Lagrangian, the
Eulerian action is
\begin{equation}\label{eq:general-potential-mfg-action}
  \int_0^T\int L(x,v_t(x))\,\dd\al_t(x)\dd t,
\end{equation}
and the same calculation yields
$v_t=-\nabla_pH(x,\nabla u_t)$.

\paragraph{Convex momentum variables.}
As in dynamic OT, write the momentum measure $m_t=\al_tv_t$.  For a density
$\rho$ and momentum density $m$, the quadratic kinetic perspective is
\begin{equation}\label{eq:mfg-perspective}
  J(\rho,m)=
  \begin{cases}
    \dfrac{\abs m^2}{2\rho},&\rho>0,\\[1mm]
    0,&(\rho,m)=(0,0),\\
    +\infty,&\rho=0,\ m\neq0.
  \end{cases}
\end{equation}
It is jointly convex and positively one-homogeneous.  Therefore
\eqref{eq:potential-mfg-primal} becomes a convex problem in $(\rho,m)$ when
$\mathcal F$ and $\mathcal G$ are convex, under the linear constraint
\begin{equation}\label{eq:mfg-linear-constraint}
  \partial_t\rho+\diverg m=\nu\Delta\rho.
\end{equation}
This is the precise bridge to Benamou--Brenier algorithms and primal--dual
splitting; see \citet{BenamouCarlier2015,BenamouCarlierSantambrogio2017}.

\begin{example}[Local congestion]
Suppose $\al=\rho\dd x$ and agents pay $F(x,\al)=f(\rho(x))$, where $f$
is increasing.  Let $U'(r)=f(r)$.  Then
\[
  \mathcal F(\al)=\int U(\rho(x))\,\dd x
\]
has first variation $F$.  The convexity of $U$ expresses aversion to
congestion.  A hard density cap is obtained by taking $U(r)=+\infty$ for
$r>\kappa$, leading to a pressure multiplier.
\end{example}

\begin{example}[Nonlocal interactions]
Let
\[
  F(x,\al)=V(x)+\int W(x,y)\,\dd\al(y),
\]
with a symmetric kernel $W(x,y)=W(y,x)$.  Then
\[
  \mathcal F(\al)=\int V\,\dd\al
  +\frac12\iint W(x,y)\,\dd\al(x)\dd\al(y).
\]
Convexity of this functional is a positive-semidefiniteness condition on the
kernel over signed measures of zero mass; it should not be confused with
pointwise convexity of $W(x,y)$.
\end{example}

\begin{remark}[Potential does not mean cooperative]
The agents remain non-cooperative.  Potentiality means that unilateral
first-order incentives integrate to one scalar functional.  Minimizing that
functional recovers the Nash equilibrium, just as a finite-dimensional
potential game can be solved by minimizing its potential.  A generic social
planner has different first-order conditions because changing the population
law changes every agent's cost.
\end{remark}

\section{JKO schemes, planning, and gradient flows}

The variational formula~\eqref{eq:potential-mfg-primal} looks close to a
Wasserstein gradient flow, but the two constructions should not be identified.

\paragraph{The JKO construction.}
For an energy $f:\Ptwo(\mathcal X)\to\RR\cup\{+\infty\}$, the JKO step with
time step $\tau>0$ is
\begin{equation}\label{eq:mfg-jko}
  \al^{k+1}\in\argmin_{\al\in\Ptwo(\mathcal X)}
  \left\{
    \frac1{2\tau}\Wass_2^2(\al^k,\al)+f(\al)
  \right\}.
\end{equation}
As $\tau\to0$, its interpolation converges under standard assumptions to
the Wasserstein gradient flow
\begin{equation}\label{eq:mfg-wgf}
  \partial_t\al_t
  =\diverg\!\left(
    \al_t\nabla\frac{\delta f}{\delta\al}(\al_t)
  \right).
\end{equation}
This evolution is initial-value and dissipative: the current measure
determines the next proximal step, and $f(\al_t)$ decreases.

By contrast, a finite-horizon MFG is an equilibrium over an entire time
interval.  Its HJB equation is terminal-value and agents anticipate future
population costs.  Even a potential MFG minimizes one global action over the
whole curve; it is closer to a problem in the calculus of variations than to a
sequence of greedy descents.

\begin{proposition}[A JKO step is a one-period planning problem]
\label{prop:jko-as-planning}
The JKO problem~\eqref{eq:mfg-jko} is equivalent to
\begin{equation}\label{eq:jko-dynamic-planning}
  \inf_{(\al_s,v_s)}
  \left\{
    \frac1{2\tau}\int_0^1\int\abs{v_s}^2\,\dd\al_s\dd s
    +f(\al_1)
  \right\},
\end{equation}
subject to
$\partial_s\al_s+\diverg(\al_sv_s)=0$ and
$\al_{s=0}=\al^k$, with $\al_1$ free.
\end{proposition}

\begin{proof}
For each candidate endpoint $\al_1$, the Benamou--Brenier formula says that
the infimum of the kinetic action over connecting curves is
$\Wass_2^2(\al^k,\al_1)$.  Minimizing also over $\al_1$ gives
\eqref{eq:mfg-jko}.
\end{proof}

The minimizing-movement construction is due to
\citet{JordanKinderlehrerOtto1998}; its metric-space theory is developed by
\citet{AmbrosioGigliSavare2008}.  Proposition~\ref{prop:jko-as-planning} is
the dynamic-transport reading of the same step.

This proposition gives a useful interpretation: a JKO step is a short-horizon
planning problem with terminal potential $f$.  Repeating the problem after
each short interval produces a receding-horizon or myopic dynamics.  A
finite-horizon MFG instead solves one rational-expectations fixed point over
$[0,T]$.

\begin{remark}[Three variational objects]
It is useful to keep three levels separate.
\begin{enumerate}[leftmargin=2em]
  \item The Benamou--Brenier problem fixes both endpoint measures and minimizes
  kinetic action.
  \item A potential MFG fixes the initial measure, usually leaves the terminal
  measure free, and adds running and terminal population potentials.
  \item A JKO step is a single short planning problem whose terminal potential
  is the energy to be decreased; iteration and the limit $\tau\to0$ create a
  gradient flow.
\end{enumerate}
The formulas overlap, but their boundary conditions and time semantics do not.
\end{remark}

\paragraph{Mean field control versus mean field games.}
Suppose a planner minimizes
\begin{equation}\label{eq:mean-field-control-cost}
  \mathbb E\!\left[
    \int_0^T \bigl(L(X_t,a_t)+F(X_t,\Law(X_t))\bigr)\,\dd t
    +G(X_T,\Law(X_T))
  \right].
\end{equation}
The planner internalizes how changing the law modifies the costs of all
agents.  Its optimality equation therefore contains derivatives of $F$ and
$G$ with respect to the measure argument.  In the MFG best response, one
agent treats $(\al_t)_t$ as fixed and does not differentiate through it.  The
two systems coincide only under special potential structures, and even then
the scalar potential minimized by the game must be distinguished from the
sum of individual costs.

\section{Monotonicity, convexity, and uniqueness}

Potentiality is sufficient for a global minimization principle, but it is not
necessary for well-posedness.  The central game-theoretic condition introduced
by Lasry and Lions is monotonicity in the measure variable.

\paragraph{Existence as a fixed point.}
The MFG construction itself suggests an existence proof: freeze a population
curve, solve the backward control problem, propagate the resulting feedback
forward, and look for a fixed point.  The following intentionally strong
version isolates this mechanism without entering the optimal regularity
theory.

\begin{theorem}[A basic parabolic existence principle]
\label{thm:mfg-basic-existence}
Let $\mathcal X=\mathbb T^d$ and $\nu>0$.  Assume that $H$ is smooth and
uniformly convex in its momentum variable, with at most quadratic growth.
Assume that $F(\cdot,\al)$ and $G(\cdot,\al)$ are uniformly bounded in
$C^2(\mathbb T^d)$ and depend continuously on $\al$ for the narrow topology.
If $\al_0$ has a smooth positive density, then the system
\eqref{eq:mfg-system} admits at least one classical solution.
\end{theorem}

\begin{proof}
Fix a continuous measure curve
$\be=(\be_t)_{0\leq t\leq T}$.  Parabolic HJB theory gives a solution
$u^\be$ with estimates controlled uniformly by the assumed bounds on $F$ and
$G$.  Hence the feedback
$b^\be=-\nabla_pH(\cdot,\nabla u^\be)$ is bounded and Hölder regular.  Solving
the linear Fokker--Planck equation with drift $b^\be$ and initial law $\al_0$
defines a new curve $\Phi(\be)$.

Parabolic compactness makes $\Phi$ map a closed convex bounded set of
Hölder-continuous density curves into a relatively compact subset of itself.
Stability of the HJB and Fokker--Planck equations with respect to their
coefficients makes $\Phi$ continuous.  Schauder's fixed-point theorem gives a
curve $\al=\Phi(\al)$.  Together with $u^\al$, it solves
\eqref{eq:mfg-system}.  The diffusion and smooth nonlocal couplings are what
make this short classical argument work; local or singular couplings require
weak-solution and compactness techniques.
\end{proof}

This fixed-point route is conceptually distinct from the direct method for a
potential game.  In the potential case, coercivity and lower
semicontinuity can produce a minimizer of~\eqref{eq:potential-mfg-primal};
the PDE system is then recovered from its optimality conditions.  General
existence theories and substantially weaker assumptions are developed in
\citet{LasryLions2007,CardaliaguetNotes2010}.

\begin{definition}[Lasry--Lions monotonicity]
A coupling $F:\mathcal X\times\Pp(\mathcal X)\to\RR$ is monotone if
\begin{equation}\label{eq:lasry-lions-monotonicity}
  \int_{\mathcal X}
  \bigl(F(x,\al)-F(x,\be)\bigr)\,\dd(\al-\be)(x)\geq0
\end{equation}
for all $\al,\be$.  It is strictly monotone if equality forces
$\al=\be$ in the relevant class.  The same definition applies to $G$.
\end{definition}

Local congestion $F(x,\rho\dd x)=f(\rho(x))$ is monotone whenever $f$ is
increasing.  A nonlocal linear coupling $F=W*\al$ is monotone exactly when
$\iint W(x,y)\dd\eta(x)\dd\eta(y)\geq0$ for every signed zero-mass measure
$\eta$.  When $F=\delta\mathcal F/\delta\al$, condition
\eqref{eq:lasry-lions-monotonicity} is the monotonicity of the first variation
and is implied by ordinary convexity of $\mathcal F$ under linear mixtures.

\begin{theorem}[Uniqueness under monotonicity]
\label{thm:mfg-uniqueness}
Let $H(x,\cdot)$ be strictly convex, and suppose $F$ and $G$ satisfy
Lasry--Lions monotonicity.  Assume two smooth solutions
$(u,\al)$ and $(\widetilde u,\widetilde\al)$ of
\eqref{eq:mfg-system} have the same initial law.  Then their equilibrium
feedbacks agree $\al_t\dd t$- and
$\widetilde\al_t\dd t$-almost everywhere.  If one coupling is strictly
monotone, the population curves agree; after a normalization, so do the value
functions.
\end{theorem}

\begin{proof}
Write $w=u-\widetilde u$ and
$\eta=\al-\widetilde\al$.  Differentiate the dual pairing
$I(t)=\int w_t\,\dd\eta_t$, use the two HJB equations and the two forward
equations, and integrate by parts.  The diffusion terms cancel.  Integrating
from $0$ to $T$, using $\eta_0=0$ and the terminal conditions, gives
\begin{align}\label{eq:mfg-uniqueness-identity}
&\int\bigl(G(x,\al_T)-G(x,\widetilde\al_T)\bigr)
     \,\dd(\al_T-\widetilde\al_T)(x)\\
&\quad+\int_0^T\!\int
  \bigl(F(x,\al_t)-F(x,\widetilde\al_t)\bigr)
  \,\dd(\al_t-\widetilde\al_t)(x)\dd t\nonumber\\
&\quad+\int_0^T\!\int
  D_H(\nabla u_t,\nabla\widetilde u_t)\,\dd\al_t\dd t
  +\int_0^T\!\int
  D_H(\nabla\widetilde u_t,\nabla u_t)
  \,\dd\widetilde\al_t\dd t=0,\nonumber
\end{align}
where
\[
  D_H(p,q)=H(x,q)-H(x,p)-\nabla_pH(x,p)\cdot(q-p)\geq0
\]
is the Bregman divergence of the convex Hamiltonian.  Every term in
\eqref{eq:mfg-uniqueness-identity} is nonnegative.  Hence all vanish.  Strict
convexity of $H$ identifies the gradients, and strict monotonicity identifies
the measures.  The PDEs then identify the remaining quantities up to the usual
additive constant in $u$ when applicable.
\end{proof}

The calculation is the MFG counterpart of a monotone-operator energy estimate.
It is not the contraction estimate of a Wasserstein gradient flow, although
both exploit a positive pairing between a difference of states and a
difference of forces.

\paragraph{Two notions of convexity.}
Lasry--Lions monotonicity concerns linear interpolation of measures and the
pairing of first variations.  Wasserstein geodesic convexity concerns
displacement interpolations.  Neither notion implies the other in complete
generality.  For a potential game, linear convexity is the natural condition
for convexity of the Eulerian momentum formulation, while displacement
convexity can provide additional estimates along transport geodesics.  The
distinction is important: importing a convergence theorem from Wasserstein
gradient flows into an MFG requires checking that the dynamics is actually a
gradient flow, not merely that the state variable is a probability measure.

\section{Canonical examples}

\subsection{Crowd motion with congestion}

Take $L(a)=\abs a^2/2$, a desired-location potential $V$, and a local
congestion cost $f(\rho)$.  The first-order system is
\begin{equation}\label{eq:mfg-crowd}
\left\{
\begin{aligned}
 -\partial_tu+\tfrac12\abs{\nabla u}^2&=V(x)+f(\rho),\\
 \partial_t\rho-\diverg(\rho\nabla u)&=0.
\end{aligned}
\right.
\end{equation}
The density anticipates bottlenecks because $u$ is solved backward.  This is
different from a purely repulsive McKean--Vlasov flow, where the drift reacts
only to the current density.  Increasing $f$ yields Lasry--Lions monotonicity;
if $f=U'$, the game is potential with internal energy $\int U(\rho)$.

\subsection{Quadratic and linear--quadratic games}

Linear dynamics, quadratic effort, and couplings through the population mean
lead to Riccati equations.  For example, let
\begin{equation}\label{eq:lq-coupling}
  F(x,\al)=\frac q2\abs{x-\bar x_\al}^2,
  \qquad
  \bar x_\al=\int y\,\dd\al(y).
\end{equation}
A quadratic ansatz for $u$ closes the HJB equation, while the forward
equation closes on the mean and covariance.  These models are the most direct
test bed for comparing finite-player Nash systems, MFG limits, and mean field
control; see \citet{HuangMalhameCaines2006,BensoussanFrehseYam2013}.

The coupling~\eqref{eq:lq-coupling} should be interpreted carefully: agents
are attracted toward the current mean when $q>0$.  Depending on signs and
terminal terms, related models can instead penalize crowding or encourage
dispersion.

\subsection{Multi-population games}

For $K$ populations, the state is a vector of measures
$\boldsymbol{\al}=(\al^1,\ldots,\al^K)$.  Each population has its own value
function and forward equation,
\begin{equation}\label{eq:multi-population-mfg}
\left\{
\begin{aligned}
 -\partial_tu^k-\nu_k\Delta u^k
   +H_k(x,\nabla u^k)&=F_k(x,\boldsymbol{\al}_t),\\
 \partial_t\al_t^k-\nu_k\Delta\al_t^k
   -\diverg\!\left(\al_t^k\nabla_pH_k(x,\nabla u_t^k)\right)&=0.
\end{aligned}
\right.
\end{equation}
Cross-couplings can encode attraction, segregation, competition for a shared
resource, or heterogeneous mobility.  Potentiality requires the vector field
of couplings $(F_1,\ldots,F_K)$ to be the first variation of one functional
of all populations, an infinite-dimensional analogue of symmetry of cross
derivatives.

\subsection{Stationary and ergodic games}

Over a long horizon, one often seeks a stationary density $\al$, a value
function $u$, and an ergodic constant $\lambda$:
\begin{equation}\label{eq:stationary-mfg}
\left\{
\begin{aligned}
 -\nu\Delta u+H(x,\nabla u)&=F(x,\al)+\lambda,\\
 -\nu\Delta\al-\diverg\!\left(\al\nabla_pH(x,\nabla u)\right)&=0,\\
 \al(\mathcal X)&=1.
\end{aligned}
\right.
\end{equation}
The scalar $\lambda$ is the equilibrium average cost per unit time.  It plays
the role of an eigenvalue in the HJB equation.  Long-time convergence of
finite-horizon systems toward ergodic equilibria is subtle because the system
is forward--backward; see \citet{CardaliaguetPorretta2019}.

\section{The master equation: a PDE on Wasserstein space}

The pair $(u_t,\al_t)$ solves one MFG for one initial distribution.  To encode
all initial laws simultaneously, introduce
\begin{equation}\label{eq:master-value}
  U(t,x,\al),
\end{equation}
the value at position $x$, time $t$, and population law $\al$.  Along an
equilibrium curve,
\begin{equation}\label{eq:master-characteristic}
  u(t,x)=U(t,x,\al_t).
\end{equation}
This is analogous to solving a finite-dimensional Hamilton--Jacobi equation
and then evaluating it along characteristics, except that the state variable
now contains a probability measure.

Write $\delta U/\delta\al(t,x,\al)(y)$ for the linear first variation in the
measure and
\begin{equation}\label{eq:lions-intrinsic-derivative}
  D_\al U(t,x,\al)(y)
  =\nabla_y\frac{\delta U}{\delta\al}(t,x,\al)(y)
\end{equation}
for its intrinsic, or Lions, derivative.  In the absence of common noise, the
formal master equation associated with~\eqref{eq:mfg-system} is
\begin{align}\label{eq:mfg-master-equation}
 -\partial_tU
 &-\nu\Delta_xU+H(x,\nabla_xU)
 -\nu\int\diverg_yD_\al U(t,x,\al)(y)\,\dd\al(y)\nonumber\\
 &+\int D_\al U(t,x,\al)(y)\cdot
   \nabla_pH\bigl(y,\nabla_xU(t,y,\al)\bigr)\,\dd\al(y)
 =F(x,\al),
\end{align}
with terminal condition $U(T,x,\al)=G(x,\al)$.

\begin{proposition}[The MFG system as characteristics of the master equation]
\label{prop:master-characteristics}
Assume $U$ is a smooth solution of~\eqref{eq:mfg-master-equation}.  Let
$\al_t$ solve the forward equation driven by
$u_t(x)=U(t,x,\al_t)$.  Then $(u,\al)$ solves the MFG
system~\eqref{eq:mfg-system}.
\end{proposition}

\begin{proof}
Differentiate~\eqref{eq:master-characteristic} along $\al_t$:
\[
  \partial_tu(t,x)=\partial_tU(t,x,\al_t)
  +\int\frac{\delta U}{\delta\al}(t,x,\al_t)(y)
       \,\partial_t\al_t(\dd y).
\]
Insert the forward Fokker--Planck equation and integrate by parts in $y$.
The diffusion gives
$\nu\int\diverg_yD_\al U\,\dd\al_t$, and the drift gives the integral
containing $D_\al U\cdot\nabla_pH$, with the signs displayed in
\eqref{eq:mfg-master-equation}.  The master equation therefore reduces to the
HJB equation along the curve.  The forward equation holds by construction.
\end{proof}

Common noise makes the population law itself stochastic and adds second-order
derivatives with respect to the measure variable.  The master equation is also
the key object for proving that the $N$-player Nash system converges to the
MFG limit; see \citet{CardaliaguetDelarueLasryLions2019}.

\begin{remark}[Measure derivatives versus Wasserstein gradients]
The scalar first variation $\delta U/\delta\al$ and the intrinsic derivative
$D_\al U=\nabla_y\delta U/\delta\al$ mirror the book's distinction between a
functional derivative and its Wasserstein gradient.  The master equation is
not itself a Wasserstein gradient-flow equation: it is a Hamilton--Jacobi
equation on a space whose state includes a measure.
\end{remark}

\section{Numerical viewpoints}

The forward--backward structure is the main numerical difficulty.  A naive
time march cannot solve both boundary conditions simultaneously.

\subsection{Forward--backward fixed-point iteration}

The simplest method alternates best response and population propagation.
Starting from a guessed curve $\al^{(0)}$, repeat
\begin{enumerate}[leftmargin=2em]
  \item solve the HJB equation backward using $\al^{(k)}$;
  \item propagate the Fokker--Planck equation forward using the resulting
  feedback, obtaining $\widehat\al^{(k+1)}$;
  \item damp the update,
  $\al^{(k+1)}=(1-\theta_k)\al^{(k)}+\theta_k\widehat\al^{(k+1)}$.
\end{enumerate}
This is a best-response or fictitious-play mechanism.  Convergence is not
automatic, but monotonicity and suitable averaging can stabilize it; see
\citet{Hadikhanloo2017,CardaliaguetHadikhanloo2017}.

\begin{center}
\fbox{\begin{minipage}{.9\textwidth}
\textbf{Forward--backward MFG iteration}\\
\textbf{Input:} initial law $\al_0$, horizon $T$, couplings $F,G$,
Hamiltonian $H$, damping sequence $(\theta_k)_k$.\\
\textbf{Initialize:} choose a population path
$\al^{(0)}=(\al_t^{(0)})_{0\leq t\leq T}$.\\
\textbf{For} $k=0,1,\ldots$\\
\hspace*{1em}$u_T^{(k)}\leftarrow G(\cdot,\al_T^{(k)})$.\\
\hspace*{1em}\textbf{Solve backward:}
$-\partial_tu^{(k)}-\nu\Delta u^{(k)}
+H(x,\nabla u^{(k)})=F(x,\al_t^{(k)})$.\\
\hspace*{1em}\textbf{Solve forward:}
$\partial_t\widehat\al^{(k+1)}-\nu\Delta\widehat\al^{(k+1)}
-\diverg(\widehat\al^{(k+1)}\nabla_pH(x,\nabla u^{(k)}))=0$,
$\widehat\al_{t=0}^{(k+1)}=\al_0$.\\
\hspace*{1em}$\al^{(k+1)}\leftarrow
(1-\theta_k)\al^{(k)}+\theta_k\widehat\al^{(k+1)}$.\\
\hspace*{1em}\textbf{Stop if} the path residual is below tolerance.\\
\textbf{Output:} approximate equilibrium $(u^{(k)},\al^{(k)})$.
\end{minipage}}
\end{center}

\subsection{Variational and primal--dual algorithms}

For potential games, the momentum formulation
\eqref{eq:mfg-perspective}--\eqref{eq:mfg-linear-constraint} is often preferable.
It exposes a convex objective under a linear space--time PDE constraint.
Augmented-Lagrangian, Douglas--Rachford, ADMM, and primal--dual methods then
alternate a local proximal step for the perspective/congestion integrand with
a global projection onto the discrete continuity or Fokker--Planck equation.
This is exactly the algorithmic architecture used for dynamic optimal
transport, with additional running and terminal potentials
\citep{BenamouCarlier2015,BenamouCarlierSantambrogio2017}.

With diffusion and quadratic Hamiltonian, a path-space entropy formulation can
turn the kinetic control cost into relative entropy with respect to Brownian
motion.  Time discretization then yields multi-marginal entropic problems and
Sinkhorn-like scaling algorithms \citep{BenamouCarlierDiMarinoNenna2019}.

\subsection{Finite differences, particles, and learning}

Monotone finite-difference schemes are natural for the HJB equation, while
mass-conservative schemes are required for the Fokker--Planck equation.
Discrete adjointness between the two is valuable: it preserves the duality
used in uniqueness and variational arguments.  Convergence of representative
schemes is developed in \citet{AchdouCamilliCapuzzoDolcetta2013}; see
\citet{AchdouLauriere2020} for a broader account.

Particle methods approximate $\al_t$ by an empirical measure.  They are
especially natural for probabilistic formulations and high-dimensional
problems, but estimating density-dependent congestion from particles requires
regularization.  Neural solvers can parametrize $u$, the feedback, or the
master field $U(t,x,\al)$; doing so does not remove the equilibrium
constraint, and training residuals must still control both the backward and
forward equations.

\section{What to remember}

\begin{enumerate}[leftmargin=2em]
  \item An MFG is the limit of a large symmetric non-cooperative game.  Its
  state is a population law, but its defining principle is Nash consistency.
  \item The value function solves a backward HJB equation; the population law
  solves a forward continuity or Fokker--Planck equation.  Their coupling is a
  fixed point between optimal response and induced law.
  \item Optimal transport supplies the kinematics, kinetic action, momentum
  perspective, and Hamilton--Jacobi duality.  Potential MFGs become
  Benamou--Brenier-type variational problems over curves of measures.
  \item A generic MFG is not a Wasserstein gradient flow.  JKO is best viewed
  here as a one-step/myopic planning construction; repeating short planning
  problems creates a dissipative flow, whereas an MFG encodes anticipation
  over a fixed horizon.
  \item Lasry--Lions monotonicity is the central uniqueness mechanism.  It is
  related to convexity of potential couplings but distinct from displacement
  convexity along Wasserstein geodesics.
  \item The master equation is a Hamilton--Jacobi equation on the space of
  positions and probability measures.  It organizes dependence on the initial
  population and the convergence from finite-player Nash systems.
\end{enumerate}

The conceptual payoff is a unified but non-reductive picture.  Optimal
transport, Wasserstein flows, mean field control, and mean field games share
the language of evolving probability measures.  They differ in who chooses
the velocity, which endpoint data are prescribed, and whether the governing
principle is transport optimality, dissipation, centralized control, or Nash
equilibrium.  Keeping these distinctions visible is what makes the common PDE
and variational machinery useful.

\bibliographystyle{plainnat}
\bibliography{references}

\end{document}
